\documentclass[10pt,a4paper]{amsart}
\usepackage[margin=19mm]{geometry}
\usepackage{amsmath,amssymb,mathtools}
\usepackage[scr=rsfs,cal=euler]{mathalfa}
\usepackage{xfrac}
\usepackage[unicode,hidelinks,bookmarks=false]{hyperref}
\newcommand{\ie}{i.e.\ }
\newcommand{\cf}{cf.\ }
\newcommand{\resp}{respectively\ }
\newcommand{\Subsection}{\subsection}
\providecommand{\nobreakdash}{\nobreak\hbox{-}}
\setlength{\emergencystretch}{2em}
\multlinegap0pt
\usepackage{amssymb}
\usepackage{xcolor}
\usepackage{enumerate}
\newtheorem{assumption}{Assumption}
\newtheorem{theorem}{Theorem}
\title{Fixed-time stabilization of systems with nilpotent matrices in the presence of large unknown delays}
\author{Kang-Kang Zhang, Emilia Fridman, \textit{Fellow, IEEE}, and Pengfei Wang}
\date{Source: arXiv:2607.24595v2 (CC BY 4.0)}
\begin{document}
\maketitle

\noindent\emph{Source and licence.} Fixed-time stabilization of systems with nilpotent matrices in the presence of large unknown delays. Version arXiv:2607.24595v2 (CC BY 4.0), distributed by arXiv under the Creative Commons Attribution 4.0 International licence, http://creativecommons.org/licenses/by/4.0/, as recorded in the arXiv licence metadata for this exact version and shown on the abstract page https://arxiv.org/abs/2607.24595v2. Full licence notice: \texttt{licenses/arxiv-2607.24595-CC-BY-4.0.txt}.\par\smallskip

\noindent\textit{Editorial scope.} Counted unit: the article's theorem vv (source0.tex lines 232--247). The article's complete original proof of it is delivered with it (source0.tex lines 248--296, one \texttt{proof} environment of 49 lines). No mathematics is written, repaired or re-derived: the statement and the proof are the article's own lines, in the article's own notation, and no retained line is substituted or re-flowed.  Retained uncounted as the source-local context the proof needs: lines 183--196; lines 198--201.  Nothing was omitted from any retained span, so the package carries no omission record.  No retained line contains a citation, so the wrapper carries no bibliography and no bibliography entry.  No retained line contains a figure environment, an image-inclusion command or drawing code, so no image is delivered and the package declares no asset dependency.  Editorial formatting only, in addition: an AMS \texttt{amsart} preamble that restates the article's own declarations and macros used by the retained lines, namely the environments \texttt{assumption}, \texttt{theorem} on separate counters, together with the standard packages those lines need.  The retained environments print in this document as Assumption 1, Theorem 1 in order of appearance, whereas the article prints the same items with the numbering of its own full document.  The complete native source of the article as distributed in the arXiv e-print for this exact version is bundled unchanged as \texttt{sources/206127/source0.tex}.
Consider the following linear continuous-time system
\begin{equation}
\left\{
\begin{array}
[c]{rl}%
\dot{x}(t) & =Ax(t)+B_{0}u(t),\\
y(t) & =x(t-\tau_{0}),
\end{array}
\right.  \label{sys_norm}%
\end{equation}
where $x(t)\in\mathbb{R}^{n}$, $A\in\mathbb{R}^{n\times n}$, $B_{0}%
\in\mathbb{R}^{n\times m}$, and $\tau_{0}\in\lbrack0,\tau_{M}]$ is an unknown
constant delay with known upper bound $\tau_{M}>0$. Here without loss of
generality, we set $y(s)=0$ for $s\leq0$.

\begin{assumption}
\label{ass0}$(A,B_{0})$ is controllable and $A$ is nilpotent with nilpotency
index $q\leq n$, namely, $A^{q}=0$ and $A^{q-1}\neq0$.
\end{assumption}

\begin{theorem}
Consider system (\ref{sys_norm}) under Assumption \ref{ass0}. Given $h\geq
\tau_{M}$, consider the smooth time-varying controller
\begin{equation}
u(t)=-K\left(  t\right)  S(t)g\left(
t\right)  y\left(  t-h\right)  ,\quad t\geq0, \label{vv}%
\end{equation}
where $l=\lfloor t/(3h)\rfloor$ and
\begin{equation}
K\left(  t\right)  =B_{0}^{\mathrm{T}}\mathrm{e}^{-A^{\mathrm{T}}\left(
t-3lh\right)  }W_{\mathrm{c}}^{-1}\mathrm{e}^{-A\left(  t-3lh-h\right)  }.
\label{K_n}%
\end{equation}
Then the closed-loop system (\ref{sys_norm}) and (\ref{vv}) satisfies the
condition that $x(t)=0$ and $u(t)=0$ for $t\geq3qh$, meaning, fixed-time stabilization.
\end{theorem}

\begin{proof}
For $t\in\lbrack3lh,(3l+2)h]$, $l\in%
%TCIMACRO{\U{2124} }%
%BeginExpansion
\mathbb{Z}
%EndExpansion
_{0}^{+}$, we have $g(t)=0$. Then the closed-loop system (\ref{sys_norm}) and
(\ref{vv}) can be represented as $\dot{x}(t)=Ax(t)$, whose solution is given
by
\begin{equation}
x(t)=\mathrm{e}^{A(t-3lh)}x(3lh). \label{11}%
\end{equation}
For $t\in((3l+2)h,(3l+3)h]$, $l\in%
%TCIMACRO{\U{2124} }%
%BeginExpansion
\mathbb{Z}
%EndExpansion
_{0}^{+}$, using $g(t)=1$, then the closed-loop system is expressed as
\begin{equation}
\dot{x}(t)=Ax(t)-B_{0}S(t)K\left(
t\right)  x\left(  t-h-\tau_{0}\right)  . \label{12}%
\end{equation}
By using the fact that $t-h-\tau_{0}\in(3lh,(3l+2)h]$, substituting (\ref{11})
into (\ref{12}), and employing the variation of constants formula, we have the
solution to (\ref{12}):
\begin{equation}
x\left(  t\right)  =\mathrm{e}^{A\left(  t-3lh\right)  }\mathit{\Delta
}(t)x\left(  3lh\right)  , \label{6}%
\end{equation}
where
\begin{align}
\mathit{\Delta}(t)=  I_{n}-\int_{2h}^{t-3lh}\mathrm{e}^{-As}B_{0}S(s) B_{0}^{\mathrm{T}}\mathrm{e}^{-A^{\mathrm{T}}s}\mathrm{d}%
sW_{\mathrm{c}}^{-1}\mathrm{e}^{-A\tau_{0}}, \label{delta}%
\end{align}
which implies that
$x((3l+3)h)  =\mathrm{e}^{3Ah}(  I_{n}-\mathrm{e}^{-A\tau_{0}})  x(
3lh)$
By iterating this, we obtain
\begin{equation}
x\left(  3hl\right)  =\left(  \mathrm{e}^{3Ah}\left(  I_{n}-\mathrm{e}%
^{-A\tau_{0}}\right)  \right)  ^{l}x\left(  0\right)  .\label{x_3lh00}%
\end{equation}
Since $A^{q}=0$, we have
$I_{n}-\mathrm{e}^{-A\tau_{0}}=A\tau_{0}-(A\tau_{0})^{2}/2!+\cdots-((-A\tau_{0})  ^{q-1})/(q-1)!,
$
which implies that $(I_{n}-\mathrm{e}^{-A\tau_{0}}) ^{q}=0$, and thus $x(3hq)=0$. From (\ref{11}) and (\ref{6}) with $l=q$, it follows $x(t)=\mathrm{e}^{A(t-3qh)}x(3qh)$ for $t\in(3qh,(3q+2)h]$ and $x(t)=\mathrm{e}^{A(t-3qh)}\mathit{\Delta}(t)x(3qh)=0$ for $t\in((3q+2)h,(3q+3)h]$ and $u(t)=-K(t)S(t)y(t-h)=-K(t)S(t)x(t-\tau_{0}-h)=0$.
Repeating the above process, for all $t\geq3qh$ we can also obtain that
$x(t)=0$ and $u(t)=0$.
\end{proof}

\end{document}
